\par\noindent\SourceLeaf{49}%
\hspace*{\parindent}The general term $r^{\text{tus}}$
\[
= \frac{1}{2^{2m-1}}\,{}^{2m}\mathfrak{B}^{(m+r)}\cos 2r\varphi\cos\varphi,
\]
\par\noindent when resolved into parts, gives
\[
\frac{1}{2^{2m-1}}\,{}^{2m}\mathfrak{B}^{(m+r)} . \frac{1}{2}\cos(2r-1)\varphi
+ \frac{1}{2^{2m-1}}\,{}^{2m}\mathfrak{B}^{(m+r)} . \frac{1}{2}\cos(2r+1)\varphi.
\]

\par\noindent Here, therefore, only odd multiples occur, and each term of the given series generates two adjacent ones. If I want term $(r-1)^{\text{tum}}$ of the new series, which is evidently multiplied by $\cos(2r-1)\varphi$, I take the latter part from term $(r-1)^{\text{to}}$ of the expanded series,
\[
= \frac{1}{2^{2m}}\,{}^{2m}\mathfrak{B}^{(m+r-1)}\cos(2r-1)\varphi,
\]
and the former part from term $r^{\text{to}}$,
\[
\frac{1}{2^{2m}}\,{}^{2m}\mathfrak{B}^{(m+r)}\cos(2r-1)\varphi
\]
and by adding them obtain
\[
\frac{1}{2^{2m}}\,{}^{2m+1}\mathfrak{B}^{(m+r)}\cos(2r-1)\varphi.
\]

\par\noindent Thus for $\cos\varphi^{2m+1}$ the general term $(r-1)^{\text{tus}}$ takes the form proposed here. Or, if you prefer, by substituting for $r$ the value $r+1$, we obtain term
$r^{\text{tus}}$
\[
\frac{1}{2^{2m}}\,{}^{2m+1}\mathfrak{B}^{(m+1+r)}\cos(2r+1)\varphi.
\]

\par\noindent It is now apparent that this law, which gives for $\cos\varphi^{2m+1}$ the term $r^{\text{tum}}$, agrees with the law assumed for $\cos\varphi^{2m-1}$ with general term $r^{\text{to}}$,
\[
= \frac{1}{2^{2m-2}}\,{}^{2m-1}\mathfrak{B}^{(m+r)}\cos(2r+1)\varphi,
\]
in such a way that the former follows immediately when we replace $m$ with $m+1$ in the latter.

\par\noindent\hspace*{\parindent}It is therefore evident, now that the indefinite progression has been established by proof, that the law is general for both even and odd powers of $\cos\varphi$.

5. Although it does not belong to the proposed investigation, I may note here that the law expressing powers of cosines by multiples of the given angle at once yields a similar law for powers of sines, merely by observing that
$\cos(90^0-\varphi) = \sin\varphi$.
